\subsubsection{Subproceso 3.A: Verificación y Control de Reglas Académicas}

\begin{enumerate}
    \item \textbf{Revisión automática de la Pre-Ficha.} \\
    \textit{Responsable: Sistema SIGA-ORCE.} \\
    El sistema revisa de forma automática la Pre-Ficha registrada para comprobar que cumpla con los prerrequisitos, el límite de créditos y la inscripción obligatoria de cursos desaprobados para alumnos en riesgo académico (Grupo Cero).

    \item \textbf{Recepción de solicitudes de rectificación.} \\
    \textit{Responsable: Dirección de Escuela / Intranet DIRCE.} \\
    Los estudiantes que tuvieron cruces de horario o falta de vacantes envían sus solicitudes para agregar, retirar o cambiar cursos a través del portal dentro de las fechas establecidas.

    \item \textbf{Evaluación y respuesta a las solicitudes de rectificación.} \\
    \textit{Responsable: Dirección de Escuela / Comisión de Matrícula.} \\
    La Dirección de Escuela revisa los pedidos presentados, aprueba o rechaza cada caso y actualiza manualmente los cursos del estudiante en el sistema.

    \item \textbf{Liberación de cupos por observaciones no corregidas.} \\
    \textit{Responsable: Sistema SIGA-ORCE.} \\
    Si la Pre-Ficha incumple alguna regla importante y el estudiante no la corrige dentro del plazo dado, el sistema anula su reserva temporal y libera las vacantes.

    \item \textbf{Emisión de la Pre-Ficha Aprobada o Subsanada.} \\
    \textit{Responsable: Sistema SIGA-ORCE.} \\
    Una vez levantadas las observaciones y aplicados los cambios aprobados, el sistema actualiza la ficha del alumno y la envía al Subproceso 3.B para su registro final.
\end{enumerate}
